% theme: sun_and_solar_system
\MajorQuestion{太陽と太陽系}
\SetRowSpaceQanda

\Instruction{次の問いに答えなさい。}
\QQ{天体が、ほかの天体のまわりを回ることを\NamedBlank{revolution}という。}{公転}
\QQ{みずから光を出してかがやく天体を\NamedBlank{star}という。}{恒星}
\QQ{\RefBlank{star}のまわりを\RefBlank{revolution}している天体を\NamedBlank{planet}という。}{惑星}
\QQ{太陽と、そのまわりを回る天体からなる集まりを\NamedBlank{solar_system}という。}{太陽系}
\QQ{太陽の表面で、周囲より温度が低く黒く見える部分を\NamedBlank{sunspot}という。}{黒点}

\Instruction{\strut}
\QQ{\RefBlank{planet}のまわりを\RefBlank{revolution}する天体を何というか。}{衛星}
\QQ{地球よりも内側を\RefBlank{revolution}する\RefBlank{planet}を何というか。}{内惑星}
\QQ{数億から数千億個の\RefBlank{star}が集まった大集団を\NamedBlank{galaxy}という。}{銀河}
\QQ{望遠鏡で太陽の像を映して観察するための板を何というか。}{太陽投影板}
\QQ{太陽の表面からふき出す、炎のような高温のガスの動きを何というか。}{プロミネンス}

\Instruction{\strut}
\QQ{地球よりも外側を\RefBlank{revolution}する\RefBlank{planet}を何というか。}{外惑星}
\QQ{直径や質量が大きく、平均密度が小さい\RefBlank{planet}を何というか。}{木星型惑星}
\QQ{海王星よりも外側を\RefBlank{revolution}する天体を何というか。}{太陽系外縁天体}
\QQ{\RefBlank{solar_system}を含む\RefBlank{galaxy}を\NamedBlank{milky_way}という。}{銀河系}
\QQ{太陽をとり巻く高温のガスの層を何というか。}{コロナ}

\Instruction{\strut}
\QQ{直径や質量が小さく、平均密度が大きい\RefBlank{planet}を何というか。}{地球型惑星}
\QQ{おもに火星と木星の公転軌道の間にあり、小形で不規則な形の天体を何というか。}{小惑星}
\QQ{おもに氷や細かいちりでできており、太陽に近づくと尾を引く天体を何というか。}{すい星}
\QQ{\RefBlank{milky_way}にある、ちりやガスの集まりを何というか。}{星雲}
\QQ{\RefBlank{sunspot}の位置が毎日少しずつ変わることから、太陽が何をしていると分かるか。}{自転}
